\section{Introduction}

Topological Integrity Gravity (TIG) introduced a structurally constrained horizon sector governed by the cubic relation

$$
x^3-x^2+\beta^3=0,
$$

where

$$
x=\frac{r_H}{2M},
\qquad
\beta=\frac{r_c}{2M}.
$$

The original TIG formulation established the existence of a critical horizon structure together with its phenomenological implications. However, the geometric assumptions underlying the cubic sector were not isolated in a fully explicit derivation framework.

The purpose of the present work is to provide a geometric reconstruction and proof of the TIG1 horizon sector under a minimal set of admissibility assumptions.

The analysis is restricted to:
- static spherical symmetry,
- Schwarzschild asymptotics,
- regular central structure,
- one intrinsic correction scale,
- and admissible horizon transition.

Within this effective geometric sector, we show that the horizon condition reduces to the TIG cubic structure.

The resulting geometry contains:
- a Schwarzschild-connected outer branch,
- an inner admissible branch,
- a degenerate critical horizon,
- and a horizonless supercritical sector.

The branch merger occurs at the critical point

$$
x_c=\frac23,
\qquad
\beta_c=\left(\frac{4}{27}\right)^{1/3},
$$

and exhibits universal saddle-node scaling behavior.

In addition to the horizon structure, we derive:
- the associated photon-sphere sector,
- the leading TIG correction to the Schwarzschild photon sphere,
- and the curvature-invariant structure of the effective geometry.

The resulting effective sector remains curvature-regular while reproducing the Schwarzschild limit asymptotically.

The scope of this paper is deliberately restricted.

The present work does not claim:
- complete nonlinear vacuum closure,
- full tensor consistency,
- a final covariant action,
- quantization,
- or a complete replacement theory for General Relativity.

These unresolved problems are deferred to the next programmatic layer, TIG3.

The goal of TIG2 is therefore the geometric reconstruction and proof of the TIG1 horizon sector within the stated effective geometric framework.
